\section*{अध्याय \ref{ch:b2:diffcalc} --- अवकल कलन}

\begin{solution}{exo:b2:diffcalc:1}
$J_f = \begin{pmatrix} 2x & -2y\\ 2y & 2x\end{pmatrix}$, $\det J_f
= 4(x^2 + y^2)$: मूल बिंदु से दूर प्रतिलोमनीय। (यह $f$ सम्मिश्र भेस में $z
\mapsto z^2$ है।)

$J_\Phi = \begin{pmatrix} \cos\theta & -r\sin\theta & 0\\
\sin\theta & r\cos\theta & 0\\ 0 & 0 & 1\end{pmatrix}$, $\det = r$: $r \neq 0$ के लिए प्रतिलोमनीय (बेलनीय निर्देशांक)।
\end{solution}

\begin{solution}{exo:b2:diffcalc:2}
$v = (a, b) \neq 0$ के लिए: $\frac{f(tv) - 0}{t} = \frac{t^3a^3}{t\cdot
t^2(a^2+b^2)} = \frac{a^3}{a^2 + b^2}$: अर्थात् हर दिशिक अवकलज विद्यमान है और उसका मान $D_v = \frac{a^3}{a^2+b^2}$ है।
परंतु $v
\mapsto D_v$ रैखिक नहीं है ($D_{(1,0)} = 1$, $D_{(0,1)} = 0$, $D_{(1,1)} = \frac12 \neq 1$): कोई भी रैखिक प्रतिचित्रण ये मान
नहीं दे सकता, अतः $f$ $0$ पर \omterm{def:b2:diffcalc:differential}{अवकलनीय} नहीं है (\omterm{def:b2:diffcalc:differential}{अवकल} को $v \mapsto D_v$ होना पड़ता)।
\end{solution}

\begin{solution}{exo:b2:diffcalc:3}
$f = x^3 + y^3 - 3xy$: क्रांतिक बिंदु $(0,0)$ और $(1,1)$ (प्रथम वर्ष का परिकलन)। हेस्सियन: $H = \begin{pmatrix} 6x & -3\\ -3 &
6y\end{pmatrix}$।
$(0,0)$ पर: अभिलक्षणिक चिह्न मिश्रित ($\det = -9 < 0$): अर्थात् काठी। $(1,1)$ पर: $\det = 27 > 0$,
अनुरेख $> 0$: अतः धनात्मक निश्चित, यथार्थ स्थानीय न्यूनतम --- और अब यह \cref{thm:b2:diffcalc:extrema}
से न्यायोचित है, केवल घोषित नहीं।

$g = x^4 + y^4 - 2(x-y)^2$: $\nabla g = (4x^3 - 4(x - y),\; 4y^3 +
4(x-y))$; क्रांतिक बिंदु $(0,0)$, $(\sqrt2, -\sqrt2)$, $(-\sqrt2, \sqrt2)$ (प्रथम वर्ष)। $(\pm\sqrt2, \mp\sqrt2)$ पर: $H =
\begin{pmatrix} 12\cdot2 - 4 & 4\\ 4 & 20\end{pmatrix} =
\begin{pmatrix} 20 & 4\\ 4 & 20\end{pmatrix}$: धनात्मक
निश्चित (विकर्ण-प्रधान; \omterm{def:b2:reduction:eigen}{अभिलक्षणिक मान} $24, 16$): अतः यथार्थ स्थानीय न्यूनतम।
और $(0,0)$ पर: $H = \begin{pmatrix} -4 & 4\\ 4 & -4\end{pmatrix}$, अपभ्रष्ट ऋणात्मक अर्धनिश्चित: जाँच चुप है; और दिशिक अध्ययन
($g(x,x) = 2x^4 > 0$, छोटा $g(x,-x) = 2x^4 - 8x^2 < 0$) काठी जैसा बिंदु दिखा देता है --- अर्थात् कोई चरम मान
नहीं।
\end{solution}

\begin{solution}{exo:b2:diffcalc:4}
$t = 1$ पर $t \mapsto f(tx)$ का अवकलन कीजिए: शृंखला नियम से $\langle \nabla f(x), x\rangle$; और समघातता से वही फलन $t^pf(x)$
है, जिसका $t = 1$ पर अवकलज $pf(x)$ है: अर्थात् ऑयलर की सर्वसमिका।

विलोम: $x \neq 0$ स्थिर कीजिए और $t > 0$ पर $\varphi(t) = f(tx) - t^p f(x)$ लीजिए। तब $\varphi'(t) = \langle\nabla f(tx), x\rangle -
pt^{p-1}f(x) = \frac1t\bigl(\langle \nabla f(tx), tx\rangle -
p\,t^pf(x)\bigr)$। और परिकल्पना ---
अर्थात् बिंदु $tx$ पर ऑयलर की सर्वसमिका --- कोष्ठक का मान $p\,f(tx) - p\,t^pf(x) =
p\,\varphi(t)$ दे देती है।
अतः $\varphi(1)
= 0$ सहित $\varphi' = \frac{p}{t}\varphi$: और इस रैखिक अवकल समीकरण का अद्वितीय हल $\varphi \equiv 0$ है (प्रथम
वर्ष की अद्वितीयता), अर्थात् $f(tx) = t^pf(x)$।
\end{solution}

\begin{solution}{exo:b2:diffcalc:5}
$f(x + h) - f(x) = \langle Ax - b, h\rangle +
\frac12\langle Ah, h\rangle$ को खोलने पर ($A$ की सममितता): सर्वत्र $\nabla f(x) = Ax -
b$ और $H_f = A$। $f$ उत्तल है तभी जब
$A$ धनात्मक अर्धनिश्चित हो (हेस्सियन जाँच, जो यहाँ वैश्विक है क्योंकि $H$
अचर है: द्वितीय-कोटि टेलर सूत्र यथार्थ है)। और यदि $A$ धनात्मक निश्चित हो:
तो अद्वितीय क्रांतिक बिंदु $x^* = A^{-1}b$ है, और $h \neq 0$ के लिए $f(x^*
+ h) - f(x^*) = \frac12\langle Ah, h\rangle \geq
\frac{\lambda_{\min}}{2}\norm h^2 > 0$: अर्थात् यथार्थ
वैश्विक न्यूनतम।
\end{solution}

\begin{solution}{exo:b2:diffcalc:6}
चूँकि $\nabla g(a) \neq 0$, कोई एक आंशिक अवकलज, मान लीजिए $\frac{\partial
g}{\partial y}(a) \neq 0$: तब अंतर्निहित फलन प्रमेय
(\cref{thm:b2:diffcalc:inverse} की साथी) $a = (a_1, a_2)$ के पास $\{g =
0\}$ का प्राचलन $y$ $C^1$, $y'(t) = -\frac{\partial_x g}{\partial_y g}(\gamma(t))$ सहित $\gamma(t) = (t, y(t))$ के रूप में
कर देती है ($g(t, y(t)) = 0$ का अवकलन कीजिए)। और एक-चर फलन $t
\mapsto f(\gamma(t))$ का $t = a_1$ पर स्थानीय चरम
मान है:
\[
0 = \frac{\dd}{\dd t}f(\gamma(t))\Big|_{a_1}
= \partial_x f(a) + \partial_y f(a)\,y'(a_1)
= \partial_x f(a) - \partial_yf(a)\frac{\partial_x
g(a)}{\partial_y g(a)} :
\]
अर्थात् सदिशों $\nabla f(a)$ और $\nabla g(a)$ के निर्देशांक अनुपाती हैं: $\lambda =
\frac{\partial_y f(a)}{\partial_y g(a)}$ सहित $\nabla f(a) = \lambda \nabla g(a)$।

अनुप्रयोग: वृत्त पर $\nabla(xy) = (y, x)$ का $(2x, 2y)$ के समांतर होना $y^2 = x^2$ को बाध्य कर देता है; और
प्रतिबंध के साथ उम्मीदवार $\pm\bigl(\tfrac{1}{\sqrt2}, \tfrac{1}{\sqrt2}\bigr)$ (मान $\frac12$) तथा $\pm\bigl(\tfrac{1}{\sqrt2},
-\tfrac{1}{\sqrt2}\bigr)$ (मान $-\frac12$) हैं: अतः महत्तम
$\frac12$, न्यूनतम $-\frac12$ (और ये प्राप्त होते हैं: वृत्त \omterm{def:b2:metric:compact}{संहत} है)।
\end{solution}

\begin{solution}{exo:b2:diffcalc:7}
$\nabla f = (2x - y + 1,\; 2y - x - 1) = 0$: हल करने पर $x =
-\frac13$, $y = \frac13$। हेस्सियन $\begin{pmatrix} 2 & -1\\ -1 &
2\end{pmatrix}$, जो धनात्मक निश्चित है: अतः
द्विघात $f$ का वैश्विक न्यूनतम, मान $f\bigl(-\frac13, \frac13\bigr) = -\frac13$।

त्रिभुज $T$ पर: अभ्यंतरीय क्रांतिक बिंदु $(-\frac13,
\frac13) \notin T$ (ऋणात्मक $x$)। भुजाएँ: $y = 0$
पर $x \in
\intcc{0}{1}$: $f = x^2 + x$, जो वर्धमान है: अतः चरम मान $0$ और $2$। $x = 0$ पर: $f = y^2 - y$, न्यूनतम
$y = \frac12$ पर $-\frac14$, और सिरों पर मान $0$ तथा $0$। $x + y = 1$ पर: $y = 1 - x$ प्रतिस्थापित
कीजिए, $f = x^2 + (1-x)^2 - x(1-x) + x - (1-x) = 3x^2 - x$; और $\intcc{0}{1}$ पर: न्यूनतम $x = \frac16$ पर $-\frac{1}{12}$, तथा मान $0$ ($x=0$ पर) और
$2$ ($x=1$ पर)। शीर्ष: $f(0,0) = 0$, $f(1,0)
= 2$, $f(0,1) = 0$। $T$ पर वैश्विक: न्यूनतम $-\frac14$
$(0,
\frac12)$ पर, महत्तम $2$ $(1, 0)$ पर।
\end{solution}

\begin{solution}{exo:b2:diffcalc:8}
$J_f = \begin{pmatrix} 1 & 2y\\ 2x & 1\end{pmatrix}$, $\det J_f =
1 - 4xy$। $0$ पर: $\det = 1 \neq 0$: अर्थात् स्थानीय अवकल समाकृतिकता (\cref{thm:b2:diffcalc:inverse}), जहाँ
\[
\dd(f^{-1})_{(0,0)} = (J_f(0))^{-1} = I_2 .
\]
किसी गोले पर प्रतिलोमनीयता: सर्वत्र $4\abs{xy} < 1$ चाहिए; और $\norm{(x,y)}_2 < r$ पर $\abs{xy} \leq \frac{x^2 + y^2}{2} <
\frac{r^2}{2}$, अतः $r = \frac{1}{\sqrt2}$ काम
करता है; और वही महत्तम है: $(x, y) = \bigl(\tfrac12, \tfrac12\bigr)$ पर, जिसका \omterm{def:b2:nvs:norm}{मानदंड} $\frac{1}{\sqrt2}$ है, $\det J_f = 0$।
\end{solution}

\begin{solution}{exo:b2:diffcalc:9}
$f(t) = (\cos t, \sin t)$: $f(0) = f(2\pi) = (1, 0)$, फिर भी $f'(t) =
(-\sin t, \cos t)$ का \omterm{def:b2:nvs:norm}{मानदंड} $1$ है, जो कभी शून्य नहीं: अर्थात् ऐसा कोई
बिंदु नहीं जहाँ अवकलज लुप्त हो --- रोल का कोई सदिश अनुरूप नहीं है। और
मध्यमान असमिका आराम से सत्य रहती है: $\norm{f(2\pi) - f(0)} = 0 \leq
2\pi \cdot \sup\norm{f'} = 2\pi$।
\end{solution}

\begin{solution}{exo:b2:diffcalc:10}
$f(x+h) - f(x) = 2\langle x, h\rangle + \norm h^2$: $\dd f_x =
2\langle x, \cdot\rangle$, $\nabla f(x) = 2x$, हेस्सियन $2I$ (अचर)। $g$ के लिए:
\[
g(x + h) - g(x) = \langle Ax, h\rangle + \langle Ah, x\rangle
+ \langle Ah, h\rangle
= \bigl\langle (A + A^{\mathsf T})x,\ h\bigr\rangle +
O(\norm h^2),
\]
अतः $\nabla g(x) = (A + A^{\mathsf T})x$ और $H_g = A +
A^{\mathsf T}$, जो अचर है। \cref{exo:b2:diffcalc:11} के अनुसार $g$ उत्तल है तभी जब $A + A^{\mathsf T}$
धनात्मक अर्धनिश्चित हो --- अर्थात् केवल $A$ का सममित भाग मायने रखता है,
और वास्तव में $g(x) =
\langle \frac{A + A^{\mathsf T}}2 x, x\rangle$।
\end{solution}

\begin{solution}{exo:b2:diffcalc:11}
$f$ उत्तल है तभी जब हर खंड तक उसका प्रतिबंधन उत्तल हो, अर्थात् तभी जब हर
$\varphi(t) = f(a + tv)$ उत्तल हो। और शृंखला नियम से $\varphi''(t) = \langle H_{a+tv}\,v,\ v\rangle$।

यदि सारे हेस्सियन धनात्मक अर्धनिश्चित हों: तो $\varphi'' \geq 0$, अतः हर $\varphi$ उत्तल है
(प्रथम वर्ष का खंड) और $f$ उत्तल। विलोमतः यदि $f$ उत्तल हो, तो हर $\varphi$
उत्तल है, अतः $\varphi''(0) \geq 0$: प्रत्येक $a$ और प्रत्येक दिशा $v$ के लिए $\langle H_a v, v\rangle \geq 0$: अर्थात्
सारे हेस्सियन धनात्मक अर्धनिश्चित हैं।
\end{solution}

\begin{solution}{exo:b2:diffcalc:12}
\begin{enumerate}
  \item $g$ पर लगाई गई मध्यमान असमिका (\cref{thm:b2:diffcalc:mvi}) से $\norm{g(x) - g(y)} \leq k\norm{x-y}$, अतः
        \[
        \norm{f(x) - f(y)} \geq \norm{x - y} - \norm{g(x) -
        g(y)} \geq (1 - k)\norm{x - y} :
        \]
        अर्थात् $f$ एकैकी है और (प्रतिबिंब पर परिभाषित) $f^{-1}$
        $\frac{1}{1-k}$-लिप्शिट्स है।
  \item $y$ स्थिर कीजिए; $T(x) = y - g(x)$ $\norm{T(x) -
        T(x')} = \norm{g(x') - g(x)} \leq k\norm{x - x'}$ पूरा करता है: अर्थात् \omterm{def:b2:metric:complete}{पूर्ण} समष्टि $\R^n$
        का संकुचन। बानाख (\cref{thm:b2:metric:banach}) अचल बिंदु $x^*
        = y - g(x^*)$ दे देती है, अर्थात् $f(x^*) = y$:
        इसलिए $f$ आच्छादक है।
  \item अतः $f$ $\R^n$ का ऐसा एकैकी आच्छादक प्रतिचित्रण है जिसका प्रतिलोम
        $\frac{1}{1-k}$-लिप्शिट्स है: यही वैश्विक प्रतिलोम प्रमेय है, जहाँ सर्वत्र
        $\dd g$ की लघुता \cref{thm:b2:diffcalc:inverse} की स्थानीय प्रतिलोमनीयता वाली परिकल्पना की जगह
        ले लेती है।
\end{enumerate}
\end{solution}

\begin{solution}{pb:b2:diffcalc:1}
\textbf{1.} $\norm{ABx} \leq \vertiii A\norm{Bx} \leq
\vertiii A\vertiii B\norm x$: $\norm x = 1$ पर उच्चतम लीजिए। आव्यूह गुणनफल, $\det$ और $\operatorname{tr}$
प्रविष्टियों के बहुपद फलन हैं, अतः \omterm{def:b2:metric:continuity}{संतत} ($\mathcal M_n(\R)
\simeq \R^{n^2}$, और सारे \omterm{def:b2:nvs:norm}{मानदंड} तुल्य: \cref{thm:b2:nvs:finitedim})।

\textbf{2.} $\sum\vertiii{X^k} \leq \sum\vertiii X^k <
\infty$: श्रेणी \omterm{def:b2:series:def}{निरपेक्ष रूप से} अभिसरित होती है, अतः अभिसरित होती
है (\cref{thm:b2:nvs:absoluteconvergence})। $(I -
X)\sum_{k\leq N}X^k = I - X^{N+1} \to I$ से: योग $(I -
X)^{-1}$ है। \omterm{def:b2:nvs:norm}{मानदंड}: $\leq \sum\vertiii X^k = \frac{1}{1 -
\vertiii X}$; और
\[
(I - X)^{-1} - I - X = \sum_{k\geq2}X^k = X^2(I - X)^{-1},
\qquad
\vertiii{X^2(I-X)^{-1}} \leq
\frac{\vertiii X^2}{1 - \vertiii X} = O(\vertiii X^2).
\]

\textbf{3.} $\vertiii{A^{-1}H}
\leq \vertiii{A^{-1}}\vertiii H < 1$ के साथ $A + H = A(I + A^{-1}H)$: प्रश्न 2 से प्रतिलोमनीय: अतः $A$ के चारों ओर
त्रिज्या $\vertiii{A^{-1}}^{-1}$ का \omterm{def:b2:metric:topology}{विवृत} गोला $GL_n(\R)$ में है। सघनता: $\det(A + \varepsilon I)$ $\varepsilon$ में घात-$n$ का बहुपद
है जिसका अग्र गुणांक $1$ है: अतः उसके परिमित मूल हैं, इसलिए ऐसे $\varepsilon_k \to 0$ हैं कि
$A + \varepsilon_kI$ प्रतिलोमनीय हो और वे $A$ पर अभिसरित हों।

\textbf{4.} $\vertiii H < \frac{1}{2\vertiii{A^{-1}}}$ के लिए:
\[
(A + H)^{-1} - A^{-1}
= \bigl[(I + A^{-1}H)^{-1} - I\bigr]A^{-1},
\]
जिसका \omterm{def:b2:nvs:norm}{मानदंड} अधिकतम $\frac{\vertiii{A^{-1}H}}{1 -
\vertiii{A^{-1}H}}\,\vertiii{A^{-1}} \leq
2\vertiii{A^{-1}}^2\vertiii H \to 0$ है: अतः $\Phi$ प्रत्येक $A \in GL_n(\R)$ पर \omterm{def:b2:metric:continuity}{संतत} है।

\textbf{5.} $(A+H)^2 = A^2 + AH + HA + H^2$: प्रतिचित्रण $H
\mapsto AH + HA$ रैखिक है और त्रुटि $H^2$ $O(\vertiii H^2)$ है। $(A + H)^k$ को
खोलकर $H$-गुणनखंडों की संख्या के अनुसार छाँटने पर: रैखिक भाग $\sum_{i=0}^{k-1}A^iHA^{k-1-i}$ है, और
$\geq 2$ गुणनखंड $H$ वाले पद $\binom k2$ गुणनफलों से परिबद्ध हैं जिनका \omterm{def:b2:nvs:norm}{मानदंड} $\leq \vertiii A^{k-2}\vertiii H^2$ के
पैमाने का है: $O(\vertiii H^2)$। योग को $kA^{k-1}H$ में समेटा नहीं जा सकता क्योंकि $H$ और
$A$ का क्रमविनिमेय होना आवश्यक नहीं --- यह योग ही सही अक्रमविनिमेय अवकलज
है।

\textbf{6.} छोटे $H$ के लिए:
\[
(A+H)^{-1} = (I + A^{-1}H)^{-1}A^{-1}
= \bigl(I - A^{-1}H + O(\vertiii H^2)\bigr)A^{-1}
= A^{-1} - A^{-1}HA^{-1} + O(\vertiii H^2) :
\]
$\dd\Phi_A(H) = -A^{-1}HA^{-1}$, जो $H$ में रैखिक है। और $n = 1$ के लिए: $\dd(1/a)(h) = -h/a^2$, अर्थात् परिचित अवकलज।

\textbf{7.} वक्र के अनुदिश शृंखला नियम: $\bigl(A(t)^{-1}\bigr)' = \dd\Phi_{A(t)}(A'(t)) =
-A(t)^{-1}A'(t)A(t)^{-1}$। $A(t) = I + tB$ पर, $t = 0$: $(I +
tB)^{-1} = I - tB + O(t^2)$।

\textbf{8.} प्रश्न 5 तथा अनुरेख की \omterm{def:b2:structures:generated}{चक्रीय} अपरिवर्त्यता से:
\[
\dd f_A(H) = \operatorname{tr}\Bigl(\sum_{i=0}^{k-1}
A^iHA^{k-1-i}\Bigr) = k\operatorname{tr}\bigl(A^{k-1}H\bigr) .
\]
और फ्रोबेनियस गुणन के विरुद्ध $\dd f_A(H) =
\operatorname{tr}\bigl((\nabla f)^{\mathsf T}H\bigr)$ के लिए $(\nabla f)^{\mathsf T} = kA^{k-1}$ आवश्यक है: $\nabla f(A) =
k\,(A^{k-1})^{\mathsf T}$।

\textbf{9.} $f(A + H) - f(A) = 2\operatorname{tr}
(A^{\mathsf T}H) + \operatorname{tr}(H^{\mathsf T}H)$: \omterm{def:b2:diffcalc:differential}{अवकल} $H \mapsto 2\operatorname{tr}(A^{\mathsf T}H) =
2\langle A, H\rangle$ है, अतः $\nabla f(A) = 2A$; और द्वितीय-कोटि पद ठीक $\norm H_F^2$ है:
अर्थात् हेस्सियन तत्समक \omterm{def:b2:quadratic:def}{द्विघात रूप} का दुगुना है, जो धनात्मक निश्चित और
अचर है, इसलिए $\norm\cdot_F^2$ यथार्थतः उत्तल है (यहाँ टेलर सूत्र यथार्थ है)।

\textbf{10.} स्तंभों में बहुरैखिकता से $\det(I + H) =
\sum_{S\subseteq\{1,\dots,n\}}\det(M_S)$, जहाँ $M_S$ का स्तंभ $j \in S$ के लिए
$h_j$ है और अन्यथा $e_j$। $S = \varnothing$ $1$ देता है; $S = \{j\}$ उस $I$ का \omterm{def:b2:linalg:det}{सारणिक} देता है जिसका
स्तंभ $j$ $h_j$ से बदल दिया गया है, अर्थात् उसकी $j$-वीं प्रविष्टि $h_{jj}$,
और योग लेने पर $\operatorname{tr}H$; और $\abs S \geq 2$ वाले हर पद में कम से कम दो स्तंभ $O(\vertiii H)$ आकार
के हैं, अतः वह $O(\vertiii H^2)$ है (परिमित-विमीय समष्टि पर बहुरैखिक प्रतिचित्रण
परिबद्ध होते हैं), और ऐसे पद परिमित हैं। अतः $\det(I + H) = 1 + \operatorname{tr}H +
O(\vertiii H^2)$: $\dd(\det)_I = \operatorname{tr}$।

\textbf{11.} $\det(A + H) = \det A\,\det(I + A^{-1}H) =
\det A\,\bigl(1 + \operatorname{tr}(A^{-1}H) +
O(\vertiii H^2)\bigr)$: \omterm{def:b2:diffcalc:differential}{अवकल} $H \mapsto
\det(A)\operatorname{tr}(A^{-1}H)$ है।

\textbf{12.} पंक्ति $i$ के अनुदिश लाप्लास प्रसार: $\det A =
\sum_j a_{ij}C_{ij}$, और सहगुणनखंड $C_{ij}$
में पंक्ति $i$ आती ही नहीं: $\frac{\partial\det}{\partial a_{ij}} = C_{ij}$। अतः
\[
\dd(\det)_A(H) = \sum_{i,j}C_{ij}h_{ij}
= \operatorname{tr}\bigl(\operatorname{com}(A)^{\mathsf T}
H\bigr)
= \operatorname{tr}\bigl(\operatorname{adj}(A)H\bigr),
\qquad
\nabla(\det)(A) = \operatorname{com}(A) .
\]
प्रतिलोमनीय $A$ के लिए $\operatorname{adj}A = \det(A)A^{-1}$ प्रश्न 11 पुनः दे देता है। और किसी $C^1$ वक्र के
अनुदिश शृंखला नियम $(\det A(t))' = \operatorname{tr}(\operatorname{adj}
(A(t))\,A'(t))$ पढ़ा जाता है: यही याकोबी सूत्र है।

\textbf{13.} $A' = MA$ और $\operatorname{adj}(A)A = \det(A)I$ के साथ:
\[
(\det A)' = \operatorname{tr}\bigl(\operatorname{adj}(A)MA
\bigr)
= \operatorname{tr}\bigl(A\operatorname{adj}(A)M\bigr)
= \det A\;\operatorname{tr}M
\]
(चक्रीयता; $A\operatorname{adj}A = \det(A) I$ भी)। और अदिश रैखिक अवकल समीकरण $y' = \operatorname{tr}(M(t))\,y$ का अद्वितीय हल $y(t) = y(0)\exp\bigl(\int_0^t
\operatorname{tr}M\bigr)$ है
(प्रथम वर्ष): यही ल्यूविल का सूत्र है।

\textbf{14.} $A \in SL_n(\R)$ पर $H = \frac1nA$ लीजिए: $\dd(\det)_A\bigl(\tfrac1nA\bigr) = \frac1n\det(A)
\operatorname{tr}(A^{-1}A) = \frac1n\cdot1\cdot n = 1 \neq 0$: अतः \omterm{def:b2:diffcalc:differential}{अवकल} कोई अशून्य रैखिक फलनिक है,
इसलिए स्तर समुच्चय के हर बिंदु पर $\R$ पर आच्छादक: अर्थात् $SL_n(\R)$ चिकना स्तर
समुच्चय है (विमा $n^2 - 1$ का)।

\textbf{15.} $\sum_k\vertiii{A^k/k!} \leq
\sum\vertiii A^k/k! = \eu^{\vertiii A}$: निरपेक्ष अभिसरण (प्रश्न 2 का पूर्णता-तर्क), और हर गोले
$\vertiii A \leq R$ पर \omterm{def:b2:funcseq:series}{सामान्य अभिसरण} ($A$ से स्वतंत्र परिबंध $R^k/k!$)। हर आंशिक योग \omterm{def:b2:metric:continuity}{संतत} है
(बहुपद); और गोलों पर एकसमान सीमा \omterm{def:b2:metric:continuity}{संतत} होती है: अतः $A \mapsto \eu^A$ \omterm{def:b2:metric:continuity}{संतत} है, जहाँ $\vertiii{\eu^A} \leq \eu^{\vertiii A}$।

\textbf{16.} दोनों श्रेणियाँ \omterm{def:b2:series:def}{निरपेक्ष रूप से} अभिसरित होती हैं, अतः \omterm{thm:b2:series:fubini}{कोशी
गुणनफल} वैध है (\cref{thm:b2:series:fubini}):
\[
\eu^A\eu^B = \sum_{n\geq0}\frac{1}{n!}\sum_{k=0}^n\binom
nkA^kB^{n-k}
= \sum_{n\geq0}\frac{(A+B)^n}{n!} = \eu^{A+B},
\]
और द्विपद सर्वसमिका को $AB = BA$ चाहिए। $B = -A$ के साथ: $\eu^A\eu^{-A} = \eu^0 = I$: अर्थात् हर $\eu^A \in GL_n(\R)$।

\textbf{17.} श्रेणी $\sum t^kA^k/k!$ और उसकी व्युत्पन्न श्रेणी $\sum t^{k-1}A^k/(k-1)! = A\sum t^{k-1}A^{k-1}/(k-1)!$ हर खंड $\abs t \leq T$ पर
\omterm{def:b2:funcseq:series}{सामान्य रूप से} अभिसरित होती हैं (परिबंध $T^k\vertiii A^k/k!$): अतः श्रेणियों के लिए अवकलन
प्रमेय (\cref{thm:b2:funcseq:seriestransfer}, प्रविष्टि-दर-प्रविष्टि लगाई गई) $\frac{\dd}{\dd t}\eu^{tA} = A\eu^{tA}$ देती है; और दाईं ओर $A$
बाहर निकालने पर $\eu^{tA}A$। इसे दोहराने पर: $C^\infty$।

\textbf{18.} $A(t) = \eu^{tA}$ $A'(t) = A\,A(t)$ पूरा करता है: अतः अचर $M = A$ के साथ ल्यूविल का सूत्र
(प्रश्न 13) $\det\eu^{tA} = \eu^{t\operatorname{tr}A}$ देता है ($t =
0$ पर मान $1$); और $t = 1$ पर $\det\eu^A = \eu^{\operatorname{tr}A}$। जाँचें: $n = 1$
स्वयं घातांकी है; शून्यंभावी $A$ के लिए $\operatorname{tr}A = 0$ और $\eu^A$ \omterm{def:b2:linalg:det}{सारणिक} $1$ वाला
एकघाती है; और $\operatorname{tr}A = 0$ $\det\eu^A = 1$ देता है: अर्थात् अनुरेख-शून्य आव्यूह $SL_n(\R)$ में
भेजे जाते हैं।

\textbf{19.} $\eu^H - I - H = \sum_{k\geq2}H^k/k!$, जिसका \omterm{def:b2:nvs:norm}{मानदंड} $\leq \vertiii H^2\eu^{\vertiii H} = O(\vertiii H^2)$ है: $\dd(\exp)_0 = \mathrm{id}$, अतः प्रतिलोमनीय। इसके
अतिरिक्त $\exp$ $C^1$ है: प्रश्न 5 से संभावित \omterm{def:b2:diffcalc:differential}{अवकल} $H \mapsto
\sum_k\frac1{k!}\sum_iA^iHA^{k-1-i}$ $A$ पर \omterm{def:b2:metric:continuity}{संतत} रूप से
निर्भर रैखिक प्रतिचित्रणों की सामान्य अभिसारी श्रेणी है (गोलों पर परिबंध
$\vertiii A^{k-1}/(k-1)!$), अतः आंशिक अवकलज विद्यमान और \omterm{def:b2:metric:continuity}{संतत} हैं (\cref{thm:b2:diffcalc:c1} तथा श्रेणी की स्थानांतरण
प्रमेय)। प्रतिलोम फलन प्रमेय (\cref{thm:b2:diffcalc:inverse}) $0$ पर लागू होती है: अर्थात् $\exp$ $0$
के किसी प्रतिवेश से $I$ के किसी प्रतिवेश पर $C^1$ अवकल समाकृतिकता है ---
$I$ के निकट आव्यूहों के लघुगणक होते हैं।

\textbf{20.} \omterm{def:b2:quadratic:adjoint}{सममित} $S = PDP^{\mathsf T}$ के लिए (वर्णक्रमीय प्रमेय): $\eu^S = P\eu^DP^{\mathsf T}$ \omterm{def:b2:quadratic:adjoint}{सममित} है और उसके
\omterm{def:b2:reduction:eigen}{अभिलक्षणिक मान} $\eu^{\lambda_i} > 0$ हैं: अतः धनात्मक निश्चित। आच्छादकता: कोई धनात्मक
निश्चित $Q =
P\operatorname{diag}(\mu_i)P^{\mathsf T}$ ($\mu_i > 0$) $S = P\operatorname{diag}(\ln\mu_i)P^{\mathsf T}$ के लिए $\eu^S$ है। एकैकीपन: $\eu^S$ अपनी अभिलक्षणिक
समष्टियाँ निर्धारित कर देता है, जो ठीक $S$ की ही हैं ($\lambda$ के लिए $S$ की
हर \omterm{def:b2:reduction:eigen}{अभिलक्षणिक समष्टि} पर $\eu^S$ $\eu^\lambda$ की तरह क्रिया करता है; और भिन्न $\lambda$
भिन्न $\eu^\lambda$ देते हैं), और अभिलक्षणिक मानों का $\ln$ लेने पर $S$ वापस मिल
जाता है। अतः $\exp$ \omterm{def:b2:quadratic:adjoint}{सममित} आव्यूहों से धनात्मक निश्चित आव्यूहों पर एकैकी
आच्छादक है।

\textbf{21.} $F(A + H) = A^{\mathsf T}A + A^{\mathsf T}H +
H^{\mathsf T}A + H^{\mathsf T}H$: $\dd F_A(H) =
A^{\mathsf T}H + H^{\mathsf T}A$ ($S_n$ में मान; त्रुटि $O(\vertiii H^2)$)। $A \in O_n$ पर और $S \in S_n$ के लिए
चुनाव $H = \frac12AS$ देता है
\[
A^{\mathsf T}\cdot\tfrac12AS + \tfrac12(AS)^{\mathsf T}A
= \tfrac12 S + \tfrac12 S^{\mathsf T} = S :
\]
अर्थात् आच्छादक। इस प्रकार $O_n = F^{-1}(I)$ विमा $n^2 - \dim S_n = \frac{n(n-1)}2$ का चिकना स्तर समुच्चय है।

\textbf{22.} $t
= 0$ पर $A(t)^{\mathsf T}A(t) = I$ का अवकलन ($A(0) = I$ के साथ): $A'(0)^{\mathsf T} + A'(0) = 0$: अर्थात् प्रतिसममित।
विलोमतः $K^{\mathsf T} = -K$ के लिए: $(\eu^{tK})^{\mathsf T}\eu^{tK} = \eu^{tK^{\mathsf T}}
\eu^{tK} = \eu^{-tK}\eu^{tK} = I$ (श्रेणी का पद-दर-पद \omterm{def:b2:linalg:transpose}{परिवर्त} लीजिए; घातांक
क्रमविनिमेय हैं): अतः वक्र $O_n$ में बना रहता है, और $t = 0$ पर उसका वेग $K$
है। इस प्रकार $I$ पर स्पर्श समष्टि $=$ प्रतिसममित आव्यूह, जिनकी विमा
अपेक्षित $\frac{n(n-1)}2$ है।

\textbf{23.} प्रतिसममित $K$ के लिए $\operatorname{tr}K = 0$, अतः $\det\eu^K = \eu^0 = 1$ (प्रश्न 18): इसलिए घातांकी
$SO_n$ में उतरता है। $n = 2$ के लिए: $J^2 = -I$, अतः $J^{2m} =
(-1)^mI$, $J^{2m+1} = (-1)^mJ$, और
\[
\eu^{\theta J}
= \Bigl(\sum_m\frac{(-1)^m\theta^{2m}}{(2m)!}\Bigr)I
+ \Bigl(\sum_m\frac{(-1)^m\theta^{2m+1}}{(2m+1)!}\Bigr)J
= \cos\theta\,I + \sin\theta\,J ,
\]
अर्थात् $\theta$ से घूर्णन: और साइन तथा कोसाइन की श्रेणियाँ \omterm{ex:b2:nvs:matrixexp}{आव्यूह घातांकी} के
भीतर बसती हैं।

\textbf{24.} प्रतिलोमनीय $A, B$ के लिए: $\operatorname{adj}(AB) =
\det(AB)(AB)^{-1} = \det(B)\det(A)B^{-1}A^{-1} =
\operatorname{adj}(B)\operatorname{adj}(A)$। सर्वसमिका के दोनों पक्ष $(A, B)$ की
प्रविष्टियों के बहुपद (अतः \omterm{def:b2:metric:continuity}{संतत}) प्रतिचित्रण हैं; और वे $\mathcal M_n\times\mathcal M_n$ के सघन
उपसमुच्चय $GL_n\times GL_n$ पर सहमत हैं (प्रश्न 3: हर गुणनखंड का सन्निकटन कीजिए), अतः
वे सर्वत्र सहमत हैं।

\textbf{25.} (क) भाग I परिमित-विमीय मानदंडित समष्टियों की पूर्णता पर चला
(निरपेक्ष अभिसारी श्रेणियाँ अभिसरित होती हैं), भाग IV उसी पर और प्रश्न 20
के लिए वर्णक्रमीय प्रमेय पर, तथा भाग V का स्थानीय लघुगणक प्रतिलोम फलन
प्रमेय पर। (ख) अवकल समीकरणों वाला अध्याय रैखिक निकायों के व्रोंस्कियन के
लिए ल्यूविल के सूत्र (प्रश्न 13) पर टिकता है, और $\frac{\dd}{\dd t}\eu^{tA} = A\eu^{tA}$ (प्रश्न 17) पर, जो
यही कहता है कि $\eu^{tA}$ $X' = AX$ को हल करता है। (ग) $\dd(\det)_I = \operatorname{tr}$ कहता है कि अनुरेख आयतन-परिवर्तन
की अत्यल्प दर है, और $\det\eu^A =
\eu^{\operatorname{tr}A}$ उसी कथन का वैश्विक समाकलन है। (घ) तृतीय वर्ष का
खंड इन संरचनाओं को नाम दे देता है: $O_n$ और $SL_n(\R)$ ली समूह हैं, $I$ पर उनकी
स्पर्श समष्टियाँ (प्रतिसममित और अनुरेख-शून्य आव्यूह) ली बीजगणित हैं, और
$\exp$ उनके बीच का पुल है।
\end{solution}
